\documentclass[11pt,a4paper]{article}
\newcommand{\gradelabel}{Grade 7}
\newcommand{\leconhdr}{Les événements et messages}
% Préambule commun aux plans de cours de robotique (Grades 7 à 12).
% Chaque plan de cours commence par \documentclass[11pt,a4paper]{article}, définit
% \gradelabel et \leconhdr (en-têtes de page), puis % Préambule commun aux plans de cours de robotique (Grades 7 à 12).
% Chaque plan de cours commence par \documentclass[11pt,a4paper]{article}, définit
% \gradelabel et \leconhdr (en-têtes de page), puis % Préambule commun aux plans de cours de robotique (Grades 7 à 12).
% Chaque plan de cours commence par \documentclass[11pt,a4paper]{article}, définit
% \gradelabel et \leconhdr (en-têtes de page), puis \input{../../preambule_robotique.tex}.
\usepackage[T1]{fontenc}
\usepackage[utf8]{inputenc}
\usepackage{lmodern}
\usepackage[french]{babel}
\usepackage[margin=2cm,headheight=14pt]{geometry}
\usepackage[most]{tcolorbox}
\usepackage{xcolor,tabularx,booktabs,array,enumitem,fancyhdr,amssymb,tikz}
\usetikzlibrary{arrows.meta,positioning}

\pagestyle{fancy}
\fancyhf{}
\lhead{\small Robotique -- \gradelabel{} -- Novembre 2026}
\rhead{\small \leconhdr}
\cfoot{\thepage}
\setlist{itemsep=1pt,topsep=2pt,leftmargin=*}
\renewcommand{\arraystretch}{1.35}
\newcolumntype{Y}{>{\raggedright\arraybackslash}X}
\newcommand{\cb}{$\square$\,}
\newcommand{\fl}{$\rightarrow$}
\newcommand{\lig}{\noindent\rule{\linewidth}{0.3pt}\\[5mm]}

\newtcolorbox{etape}[3]{enhanced,breakable,colback=#1!6,colframe=#1!70!black,
  coltitle=white,fonttitle=\bfseries,title={#2\hfill #3},boxrule=0.8pt,arc=2mm,
  before skip=8pt,after skip=8pt}
\newtcolorbox{astuce}[1][Différenciation et besoins particuliers]{breakable,
  colback=yellow!12,colframe=yellow!50!black,fonttitle=\bfseries\small,
  title={#1},boxrule=0.5pt,arc=1mm,before skip=6pt,after skip=6pt}
\newtcolorbox{code}[1][Code]{breakable,colback=black!4,colframe=black!60,
  fonttitle=\bfseries\small,title={#1},boxrule=0.5pt,arc=1mm,fontupper=\ttfamily\small,
  before skip=6pt,after skip=6pt}
% Blocs de code coloré (écrire le code sans accents : comme dans l'éditeur)
\newtcblisting{arduino}[1][Code Arduino]{listing only,breakable,colback=black!4,
  colframe=black!60,fonttitle=\bfseries\small,title={#1},boxrule=0.5pt,arc=1mm,
  before skip=6pt,after skip=6pt,
  listing options={language=C++,basicstyle=\ttfamily\small,columns=fullflexible,
  keepspaces=true,showstringspaces=false,commentstyle=\color{green!40!black}\itshape,
  keywordstyle=\color{blue!70!black}\bfseries}}
\newtcblisting{python}[1][Code Python]{listing only,breakable,colback=black!4,
  colframe=black!60,fonttitle=\bfseries\small,title={#1},boxrule=0.5pt,arc=1mm,
  before skip=6pt,after skip=6pt,
  listing options={language=Python,basicstyle=\ttfamily\small,columns=fullflexible,
  keepspaces=true,showstringspaces=false,commentstyle=\color{green!40!black}\itshape,
  keywordstyle=\color{blue!70!black}\bfseries}}
\newcommand{\entete}[2]{%
  \begin{tcolorbox}[colback=#1!10,colframe=#1!70!black,boxrule=0.8pt,arc=2mm]
  {\Large\bfseries #2}\hfill Nom : \rule{4cm}{0.3pt}\quad Date : \rule{2.2cm}{0.3pt}
  \end{tcolorbox}}

% Bandeau de titre : #1 titre, #2 semaine
\newcommand{\titreblock}[2]{%
\begin{tcolorbox}[colback=teal!10,colframe=teal!70!black,arc=2mm,boxrule=1pt]
\centering
{\LARGE\bfseries #1}\\[2mm]
{\large Robotique -- \gradelabel{} -- Plan de cours de 2 heures}\\[1mm]
Enseignant : Francis Ki \quad|\quad #2
\end{tcolorbox}}

% Tableau du minutage : 6 contenus (accroche, mini-leçon, guidée, autonome, mise en commun, exit ticket)
\newcommand{\minutage}[6]{%
\begin{tabularx}{\linewidth}{>{\bfseries}p{2.2cm}p{2.6cm}Y p{3cm}}
\toprule
Minutes & Étape & Contenu & Organisation\\
\midrule
0 -- 10 & Accroche & #1 & Classe entière\\
10 -- 30 & Mini-leçon & #2 & Classe entière\\
30 -- 60 & Pratique guidée & #3 & Collectif puis binômes\\
60 -- 95 & Pratique autonome & #4 & Individuel sur poste\\
95 -- 110 & Mise en commun & #5 & Classe entière\\
110 -- 120 & Exit ticket & #6 & Individuel sur poste\\
\bottomrule
\end{tabularx}}

% Étayage progressif : #1 aide 1, #2 aide 2, #3 aide 3
\newcommand{\aides}[3]{%
\textbf{Étayage progressif de l'aide.}
\begin{center}
\begin{tabularx}{\linewidth}{>{\bfseries}p{2.4cm}Y p{4.4cm}}
\toprule
Niveau d'aide & Contenu & Qui intervient\\
\midrule
Aide 1 & #1 & Le binôme seul\\
Aide 2 & #2 & Le binôme avec la carte d'aide\\
Aide 3 & #3 & L'enseignant\\
\bottomrule
\end{tabularx}
\end{center}}

% Les trois groupes de la pratique autonome : #1 tâche A, #2 tâche B, #3 tâche C
\newcommand{\groupes}[3]{%
\textbf{Consigne commune\,:} chaque élève ouvre le fichier de son groupe, le renomme avec son prénom et travaille seul sur son poste. L'enseignant passe d'abord auprès du groupe A pendant 10 min.

\medskip
\begin{tabularx}{\linewidth}{>{\bfseries}p{3cm}Y p{2.2cm}}
\toprule
Groupe & Tâche sur poste & Annexe\\
\midrule
A : en dessous du niveau & #1 & Annexe C\\
B : au niveau & #2 & Annexe D\\
C : au-dessus du niveau & #3 & Annexe E\\
\bottomrule
\end{tabularx}}

% Mise en commun : #1 groupe A, #2 groupe B, #3 groupe C, #4 conclusion
\newcommand{\miseencommun}[4]{%
\begin{etape}{teal}{5. Mise en commun}{95 -- 110 min}
\begin{enumerate}
 \item (2 min) Retour au calme, écrans mis en veille.
 \item (9 min) Trois passages courts de 3 min, avec valorisation :
 \begin{itemize}
  \item \textbf{Groupe A} : #1
  \item \textbf{Groupe B} : #2
  \item \textbf{Groupe C} : #3
 \end{itemize}
 \item (4 min) Conclusion : #4
\end{enumerate}
\end{etape}}

% Clôture : #1 question 1, #2 question 2
\newcommand{\cloture}[2]{%
\begin{etape}{gray}{6. Clôture : exit ticket}{110 -- 120 min}
Formulaire en ligne ou papier (annexe G) :
\begin{enumerate}
 \item #1
 \item #2
 \item Auto-évaluation : \emph{Je maîtrise / Je progresse / J'ai besoin d'aide.}
\end{enumerate}
Lire les réponses avant la séance suivante pour ajuster les groupes de besoin.
\end{etape}}

% Grille d'observation : #1 #2 #3 intitulés de colonnes
\newcommand{\grille}[3]{%
\subsection*{Grille d'observation (enseignant)}
\begin{tabularx}{\linewidth}{Y p{2.4cm} p{2.4cm} p{2.4cm} Y}
\toprule
Élève & #1 & #2 & #3 & Groupe suivant\\
\midrule
 & & & & \\ & & & & \\ & & & & \\ & & & & \\ & & & & \\
\bottomrule
\end{tabularx}}

% Signature de l'exit ticket
\newcommand{\autoeval}{Je me sens : \cb Je maîtrise \quad \cb Je progresse \quad \cb J'ai besoin d'aide}
.
\usepackage[T1]{fontenc}
\usepackage[utf8]{inputenc}
\usepackage{lmodern}
\usepackage[french]{babel}
\usepackage[margin=2cm,headheight=14pt]{geometry}
\usepackage[most]{tcolorbox}
\usepackage{xcolor,tabularx,booktabs,array,enumitem,fancyhdr,amssymb,tikz}
\usetikzlibrary{arrows.meta,positioning}

\pagestyle{fancy}
\fancyhf{}
\lhead{\small Robotique -- \gradelabel{} -- Novembre 2026}
\rhead{\small \leconhdr}
\cfoot{\thepage}
\setlist{itemsep=1pt,topsep=2pt,leftmargin=*}
\renewcommand{\arraystretch}{1.35}
\newcolumntype{Y}{>{\raggedright\arraybackslash}X}
\newcommand{\cb}{$\square$\,}
\newcommand{\fl}{$\rightarrow$}
\newcommand{\lig}{\noindent\rule{\linewidth}{0.3pt}\\[5mm]}

\newtcolorbox{etape}[3]{enhanced,breakable,colback=#1!6,colframe=#1!70!black,
  coltitle=white,fonttitle=\bfseries,title={#2\hfill #3},boxrule=0.8pt,arc=2mm,
  before skip=8pt,after skip=8pt}
\newtcolorbox{astuce}[1][Différenciation et besoins particuliers]{breakable,
  colback=yellow!12,colframe=yellow!50!black,fonttitle=\bfseries\small,
  title={#1},boxrule=0.5pt,arc=1mm,before skip=6pt,after skip=6pt}
\newtcolorbox{code}[1][Code]{breakable,colback=black!4,colframe=black!60,
  fonttitle=\bfseries\small,title={#1},boxrule=0.5pt,arc=1mm,fontupper=\ttfamily\small,
  before skip=6pt,after skip=6pt}
% Blocs de code coloré (écrire le code sans accents : comme dans l'éditeur)
\newtcblisting{arduino}[1][Code Arduino]{listing only,breakable,colback=black!4,
  colframe=black!60,fonttitle=\bfseries\small,title={#1},boxrule=0.5pt,arc=1mm,
  before skip=6pt,after skip=6pt,
  listing options={language=C++,basicstyle=\ttfamily\small,columns=fullflexible,
  keepspaces=true,showstringspaces=false,commentstyle=\color{green!40!black}\itshape,
  keywordstyle=\color{blue!70!black}\bfseries}}
\newtcblisting{python}[1][Code Python]{listing only,breakable,colback=black!4,
  colframe=black!60,fonttitle=\bfseries\small,title={#1},boxrule=0.5pt,arc=1mm,
  before skip=6pt,after skip=6pt,
  listing options={language=Python,basicstyle=\ttfamily\small,columns=fullflexible,
  keepspaces=true,showstringspaces=false,commentstyle=\color{green!40!black}\itshape,
  keywordstyle=\color{blue!70!black}\bfseries}}
\newcommand{\entete}[2]{%
  \begin{tcolorbox}[colback=#1!10,colframe=#1!70!black,boxrule=0.8pt,arc=2mm]
  {\Large\bfseries #2}\hfill Nom : \rule{4cm}{0.3pt}\quad Date : \rule{2.2cm}{0.3pt}
  \end{tcolorbox}}

% Bandeau de titre : #1 titre, #2 semaine
\newcommand{\titreblock}[2]{%
\begin{tcolorbox}[colback=teal!10,colframe=teal!70!black,arc=2mm,boxrule=1pt]
\centering
{\LARGE\bfseries #1}\\[2mm]
{\large Robotique -- \gradelabel{} -- Plan de cours de 2 heures}\\[1mm]
Enseignant : Francis Ki \quad|\quad #2
\end{tcolorbox}}

% Tableau du minutage : 6 contenus (accroche, mini-leçon, guidée, autonome, mise en commun, exit ticket)
\newcommand{\minutage}[6]{%
\begin{tabularx}{\linewidth}{>{\bfseries}p{2.2cm}p{2.6cm}Y p{3cm}}
\toprule
Minutes & Étape & Contenu & Organisation\\
\midrule
0 -- 10 & Accroche & #1 & Classe entière\\
10 -- 30 & Mini-leçon & #2 & Classe entière\\
30 -- 60 & Pratique guidée & #3 & Collectif puis binômes\\
60 -- 95 & Pratique autonome & #4 & Individuel sur poste\\
95 -- 110 & Mise en commun & #5 & Classe entière\\
110 -- 120 & Exit ticket & #6 & Individuel sur poste\\
\bottomrule
\end{tabularx}}

% Étayage progressif : #1 aide 1, #2 aide 2, #3 aide 3
\newcommand{\aides}[3]{%
\textbf{Étayage progressif de l'aide.}
\begin{center}
\begin{tabularx}{\linewidth}{>{\bfseries}p{2.4cm}Y p{4.4cm}}
\toprule
Niveau d'aide & Contenu & Qui intervient\\
\midrule
Aide 1 & #1 & Le binôme seul\\
Aide 2 & #2 & Le binôme avec la carte d'aide\\
Aide 3 & #3 & L'enseignant\\
\bottomrule
\end{tabularx}
\end{center}}

% Les trois groupes de la pratique autonome : #1 tâche A, #2 tâche B, #3 tâche C
\newcommand{\groupes}[3]{%
\textbf{Consigne commune\,:} chaque élève ouvre le fichier de son groupe, le renomme avec son prénom et travaille seul sur son poste. L'enseignant passe d'abord auprès du groupe A pendant 10 min.

\medskip
\begin{tabularx}{\linewidth}{>{\bfseries}p{3cm}Y p{2.2cm}}
\toprule
Groupe & Tâche sur poste & Annexe\\
\midrule
A : en dessous du niveau & #1 & Annexe C\\
B : au niveau & #2 & Annexe D\\
C : au-dessus du niveau & #3 & Annexe E\\
\bottomrule
\end{tabularx}}

% Mise en commun : #1 groupe A, #2 groupe B, #3 groupe C, #4 conclusion
\newcommand{\miseencommun}[4]{%
\begin{etape}{teal}{5. Mise en commun}{95 -- 110 min}
\begin{enumerate}
 \item (2 min) Retour au calme, écrans mis en veille.
 \item (9 min) Trois passages courts de 3 min, avec valorisation :
 \begin{itemize}
  \item \textbf{Groupe A} : #1
  \item \textbf{Groupe B} : #2
  \item \textbf{Groupe C} : #3
 \end{itemize}
 \item (4 min) Conclusion : #4
\end{enumerate}
\end{etape}}

% Clôture : #1 question 1, #2 question 2
\newcommand{\cloture}[2]{%
\begin{etape}{gray}{6. Clôture : exit ticket}{110 -- 120 min}
Formulaire en ligne ou papier (annexe G) :
\begin{enumerate}
 \item #1
 \item #2
 \item Auto-évaluation : \emph{Je maîtrise / Je progresse / J'ai besoin d'aide.}
\end{enumerate}
Lire les réponses avant la séance suivante pour ajuster les groupes de besoin.
\end{etape}}

% Grille d'observation : #1 #2 #3 intitulés de colonnes
\newcommand{\grille}[3]{%
\subsection*{Grille d'observation (enseignant)}
\begin{tabularx}{\linewidth}{Y p{2.4cm} p{2.4cm} p{2.4cm} Y}
\toprule
Élève & #1 & #2 & #3 & Groupe suivant\\
\midrule
 & & & & \\ & & & & \\ & & & & \\ & & & & \\ & & & & \\
\bottomrule
\end{tabularx}}

% Signature de l'exit ticket
\newcommand{\autoeval}{Je me sens : \cb Je maîtrise \quad \cb Je progresse \quad \cb J'ai besoin d'aide}
.
\usepackage[T1]{fontenc}
\usepackage[utf8]{inputenc}
\usepackage{lmodern}
\usepackage[french]{babel}
\usepackage[margin=2cm,headheight=14pt]{geometry}
\usepackage[most]{tcolorbox}
\usepackage{xcolor,tabularx,booktabs,array,enumitem,fancyhdr,amssymb,tikz}
\usetikzlibrary{arrows.meta,positioning}

\pagestyle{fancy}
\fancyhf{}
\lhead{\small Robotique -- \gradelabel{} -- Novembre 2026}
\rhead{\small \leconhdr}
\cfoot{\thepage}
\setlist{itemsep=1pt,topsep=2pt,leftmargin=*}
\renewcommand{\arraystretch}{1.35}
\newcolumntype{Y}{>{\raggedright\arraybackslash}X}
\newcommand{\cb}{$\square$\,}
\newcommand{\fl}{$\rightarrow$}
\newcommand{\lig}{\noindent\rule{\linewidth}{0.3pt}\\[5mm]}

\newtcolorbox{etape}[3]{enhanced,breakable,colback=#1!6,colframe=#1!70!black,
  coltitle=white,fonttitle=\bfseries,title={#2\hfill #3},boxrule=0.8pt,arc=2mm,
  before skip=8pt,after skip=8pt}
\newtcolorbox{astuce}[1][Différenciation et besoins particuliers]{breakable,
  colback=yellow!12,colframe=yellow!50!black,fonttitle=\bfseries\small,
  title={#1},boxrule=0.5pt,arc=1mm,before skip=6pt,after skip=6pt}
\newtcolorbox{code}[1][Code]{breakable,colback=black!4,colframe=black!60,
  fonttitle=\bfseries\small,title={#1},boxrule=0.5pt,arc=1mm,fontupper=\ttfamily\small,
  before skip=6pt,after skip=6pt}
% Blocs de code coloré (écrire le code sans accents : comme dans l'éditeur)
\newtcblisting{arduino}[1][Code Arduino]{listing only,breakable,colback=black!4,
  colframe=black!60,fonttitle=\bfseries\small,title={#1},boxrule=0.5pt,arc=1mm,
  before skip=6pt,after skip=6pt,
  listing options={language=C++,basicstyle=\ttfamily\small,columns=fullflexible,
  keepspaces=true,showstringspaces=false,commentstyle=\color{green!40!black}\itshape,
  keywordstyle=\color{blue!70!black}\bfseries}}
\newtcblisting{python}[1][Code Python]{listing only,breakable,colback=black!4,
  colframe=black!60,fonttitle=\bfseries\small,title={#1},boxrule=0.5pt,arc=1mm,
  before skip=6pt,after skip=6pt,
  listing options={language=Python,basicstyle=\ttfamily\small,columns=fullflexible,
  keepspaces=true,showstringspaces=false,commentstyle=\color{green!40!black}\itshape,
  keywordstyle=\color{blue!70!black}\bfseries}}
\newcommand{\entete}[2]{%
  \begin{tcolorbox}[colback=#1!10,colframe=#1!70!black,boxrule=0.8pt,arc=2mm]
  {\Large\bfseries #2}\hfill Nom : \rule{4cm}{0.3pt}\quad Date : \rule{2.2cm}{0.3pt}
  \end{tcolorbox}}

% Bandeau de titre : #1 titre, #2 semaine
\newcommand{\titreblock}[2]{%
\begin{tcolorbox}[colback=teal!10,colframe=teal!70!black,arc=2mm,boxrule=1pt]
\centering
{\LARGE\bfseries #1}\\[2mm]
{\large Robotique -- \gradelabel{} -- Plan de cours de 2 heures}\\[1mm]
Enseignant : Francis Ki \quad|\quad #2
\end{tcolorbox}}

% Tableau du minutage : 6 contenus (accroche, mini-leçon, guidée, autonome, mise en commun, exit ticket)
\newcommand{\minutage}[6]{%
\begin{tabularx}{\linewidth}{>{\bfseries}p{2.2cm}p{2.6cm}Y p{3cm}}
\toprule
Minutes & Étape & Contenu & Organisation\\
\midrule
0 -- 10 & Accroche & #1 & Classe entière\\
10 -- 30 & Mini-leçon & #2 & Classe entière\\
30 -- 60 & Pratique guidée & #3 & Collectif puis binômes\\
60 -- 95 & Pratique autonome & #4 & Individuel sur poste\\
95 -- 110 & Mise en commun & #5 & Classe entière\\
110 -- 120 & Exit ticket & #6 & Individuel sur poste\\
\bottomrule
\end{tabularx}}

% Étayage progressif : #1 aide 1, #2 aide 2, #3 aide 3
\newcommand{\aides}[3]{%
\textbf{Étayage progressif de l'aide.}
\begin{center}
\begin{tabularx}{\linewidth}{>{\bfseries}p{2.4cm}Y p{4.4cm}}
\toprule
Niveau d'aide & Contenu & Qui intervient\\
\midrule
Aide 1 & #1 & Le binôme seul\\
Aide 2 & #2 & Le binôme avec la carte d'aide\\
Aide 3 & #3 & L'enseignant\\
\bottomrule
\end{tabularx}
\end{center}}

% Les trois groupes de la pratique autonome : #1 tâche A, #2 tâche B, #3 tâche C
\newcommand{\groupes}[3]{%
\textbf{Consigne commune\,:} chaque élève ouvre le fichier de son groupe, le renomme avec son prénom et travaille seul sur son poste. L'enseignant passe d'abord auprès du groupe A pendant 10 min.

\medskip
\begin{tabularx}{\linewidth}{>{\bfseries}p{3cm}Y p{2.2cm}}
\toprule
Groupe & Tâche sur poste & Annexe\\
\midrule
A : en dessous du niveau & #1 & Annexe C\\
B : au niveau & #2 & Annexe D\\
C : au-dessus du niveau & #3 & Annexe E\\
\bottomrule
\end{tabularx}}

% Mise en commun : #1 groupe A, #2 groupe B, #3 groupe C, #4 conclusion
\newcommand{\miseencommun}[4]{%
\begin{etape}{teal}{5. Mise en commun}{95 -- 110 min}
\begin{enumerate}
 \item (2 min) Retour au calme, écrans mis en veille.
 \item (9 min) Trois passages courts de 3 min, avec valorisation :
 \begin{itemize}
  \item \textbf{Groupe A} : #1
  \item \textbf{Groupe B} : #2
  \item \textbf{Groupe C} : #3
 \end{itemize}
 \item (4 min) Conclusion : #4
\end{enumerate}
\end{etape}}

% Clôture : #1 question 1, #2 question 2
\newcommand{\cloture}[2]{%
\begin{etape}{gray}{6. Clôture : exit ticket}{110 -- 120 min}
Formulaire en ligne ou papier (annexe G) :
\begin{enumerate}
 \item #1
 \item #2
 \item Auto-évaluation : \emph{Je maîtrise / Je progresse / J'ai besoin d'aide.}
\end{enumerate}
Lire les réponses avant la séance suivante pour ajuster les groupes de besoin.
\end{etape}}

% Grille d'observation : #1 #2 #3 intitulés de colonnes
\newcommand{\grille}[3]{%
\subsection*{Grille d'observation (enseignant)}
\begin{tabularx}{\linewidth}{Y p{2.4cm} p{2.4cm} p{2.4cm} Y}
\toprule
Élève & #1 & #2 & #3 & Groupe suivant\\
\midrule
 & & & & \\ & & & & \\ & & & & \\ & & & & \\ & & & & \\
\bottomrule
\end{tabularx}}

% Signature de l'exit ticket
\newcommand{\autoeval}{Je me sens : \cb Je maîtrise \quad \cb Je progresse \quad \cb J'ai besoin d'aide}


\begin{document}
\titreblock{Les événements et messages}{Semaine du 16 au 20 novembre 2026}

%==================================================================
\section{Informations générales}
\begin{tabularx}{\linewidth}{>{\bfseries}p{3.6cm}Y}
\toprule
Titre & Les événements et messages (Scratch)\\
Classe / effectif & Grade 7, 11 à 20 élèves, niveaux hétérogènes, élèves à besoins particuliers\\
Durée & 2 heures (120 min)\\
Lieu & Salle informatique, un poste par élève\\
Matériel et logiciels &
\begin{itemize}
 \item Un poste par élève avec Scratch 3 (éditeur en ligne ou hors ligne).
 \item Vidéoprojecteur ou tableau interactif.
 \item Projets de départ \texttt{.sb3} pour les trois niveaux, dans un dossier partagé.
\end{itemize}\\
Pré-requis &
Boucles et conditions (deux séances précédentes). Savoir ajouter un lutin et un arrière-plan.\\
\bottomrule
\end{tabularx}

\begin{astuce}[Préparation avant le cours (environ 40 min)]
\begin{enumerate}
 \item Préparer trois projets : A (deux lutins qui dialoguent, blocs d'envoi et de réception manquants), B (histoire en trois scènes sans messages), C (jeu à deux lutins sans communication).
 \item Tester les projets sur les postes élèves.
 \item Créer le formulaire de sortie (annexe G) ou imprimer l'annexe.
 \item Reconstituer les groupes à partir de l'exit ticket de la séance précédente.
\end{enumerate}
\end{astuce}

%==================================================================
\section{Objectifs d'apprentissage}
\begin{tcolorbox}[colback=gray!8,colframe=gray!60!black,arc=1mm,boxrule=0.5pt]
\textbf{Objectifs du programme}
\begin{itemize}
 \item Déclencher une action par événement.
 \item Faire communiquer plusieurs personnages Scratch.
\end{itemize}
\textbf{À la fin de la séance, l'élève est capable\,:}
\begin{itemize}
 \item d'utiliser les blocs d'événement (\texttt{quand drapeau cliqué}, \texttt{quand touche pressée}, \texttt{quand ce lutin est cliqué})\,;
 \item d'envoyer et de recevoir un message entre deux lutins\,;
 \item d'expliquer l'ordre dans lequel les actions se produisent.
\end{itemize}
\textbf{Compétences numériques\,:} gérer plusieurs lutins, nommer ses messages, tester et corriger.
\end{tcolorbox}

%==================================================================
\section{Déroulé minuté (120 min)}
\minutage{Le relais de messages}{Événement, envoi et réception}{Un dialogue entre deux lutins, puis un changement de décor en binômes}{Trois niveaux de besoin}{Valoriser chaque groupe}{Trois questions sur poste}

%------------------------------------------------------------------
\begin{etape}{orange}{1. Accroche : le relais de messages}{0 -- 10 min}
\textbf{Écrans éteints.}
\begin{enumerate}
 \item (6 min) Trois élèves en ligne. L'enseignant tape sur l'épaule du premier, qui dit «~Bonjour~» et tape sur l'épaule du deuxième, qui répond «~Salut~» puis tape sur l'épaule du troisième, qui lève la main. La classe observe qui agit quand.
 \item (4 min) Question-problème : \textbf{«~Comment un lutin peut-il dire à un autre lutin de commencer ?~»}
\end{enumerate}
\end{etape}

%------------------------------------------------------------------
\begin{etape}{blue}{2. Mini-leçon : événements et messages}{10 -- 30 min}
\textbf{Support\,:} annexe A. Quatre temps de 5 min.

\textbf{Temps 1 : un événement.} Une action programmée démarre quand \emph{quelque chose se produit} : drapeau vert, touche pressée, clic sur le lutin.

\textbf{Temps 2 : les messages.} Un lutin peut \emph{envoyer} un message. Un autre lutin le \emph{reçoit} et démarre un script.

\begin{center}
\begin{tikzpicture}[box/.style={draw,rounded corners,thick,minimum height=1cm,text width=3.2cm,align=center,font=\small},>={Stealth[length=3mm]},node distance=1.2cm]
\node[box,fill=orange!20] (a) {Lutin 1\\ dire «~Bonjour~»\\ \texttt{envoyer salut}};
\node[box,fill=blue!15,right=of a] (b) {Message\\ \texttt{salut}};
\node[box,fill=green!15,right=of b] (c) {Lutin 2\\ \texttt{quand je reçois salut}\\ dire «~Salut !~»};
\draw[->,thick] (a)--(b); \draw[->,thick] (b)--(c);
\end{tikzpicture}
\end{center}

\textbf{Temps 3 : l'ordre des actions.} Quand plusieurs scripts démarrent en même temps, ils s'exécutent \emph{en parallèle}. Pour les enchaîner : \texttt{envoyer à tous et attendre}.

\textbf{Temps 4 : bien nommer.} Choisir des noms clairs pour les messages : \texttt{debut\_scene2} plutôt que \texttt{message1}.

\begin{astuce}
Utiliser des couleurs par lutin dans les schémas. Donner un modèle de nom de message. Lire à voix haute : «~quand je reçois salut, alors\ldots~»
\end{astuce}
\end{etape}

%------------------------------------------------------------------
\begin{etape}{green}{3. Pratique guidée}{30 -- 60 min}
\textbf{Fiche\,:} annexe B.

\textbf{3a. Exemple collectif sur poste (12 min).} Chat et chien : le chat dit «~Bonjour~» quand on clique dessus, puis envoie le message \texttt{salut}. Le chien reçoit \texttt{salut} et répond.

\textbf{3b. Second exemple en binômes (18 min).} Un clic sur un bouton envoie le message \texttt{nuit}. L'arrière-plan change et un lutin dit «~Bonne nuit~». Alternance des rôles à mi-temps.

\aides{Le projet du 3a et le schéma de l'annexe A}{Questions guides : \emph{Qui envoie ? Qui reçoit ? Quel nom donnes-tu au message ?}}{L'enseignant crée le message avec le binôme, qui termine le script}

\begin{astuce}
Proposer des blocs déjà posés dans le projet pour les élèves lents. Autoriser à tester lutin par lutin avant de lancer l'ensemble.
\end{astuce}
\end{etape}

%------------------------------------------------------------------
\begin{etape}{violet}{4. Pratique autonome différenciée}{60 -- 95 min}
\groupes{Projet à compléter : un dialogue entre deux lutins dont les blocs \texttt{envoyer} et \texttt{quand je reçois} sont à relier avec des flèches de couleur, puis ajouter une réplique.}{Créer une histoire en trois scènes (arrière-plans différents) avec trois lutins qui s'envoient des messages pour passer d'une scène à l'autre.}{Programmer un jeu à deux lutins qui communiquent : un message \texttt{touche} quand le joueur est atteint, un message \texttt{fin} qui arrête les scripts, et une variable \texttt{vies}. Documenter les messages dans un tableau.}

\medskip
\textbf{Plan B si un élève termine tôt\,:} ajouter un bouton «~Rejouer~» qui envoie un message à tous les lutins.
\end{etape}

%------------------------------------------------------------------
\miseencommun{un élève montre le dialogue et nomme qui envoie et qui reçoit. Valoriser la précision du vocabulaire.}{un binôme présente son histoire et enchaîne les trois scènes. Valoriser la fluidité.}{un élève présente son jeu et le tableau des messages. Valoriser la structure.}{\textit{un événement déclenche une action et un message permet à deux lutins de se parler.}}

\cloture{Quel bloc permet à un lutin de réagir quand un autre lutin lui envoie un signal ?}{Pourquoi est-il utile de donner un nom clair à un message ?}

%==================================================================
% ANNEXES
%==================================================================
\newpage
\section*{Annexe A : Support résumé de la mini-leçon}
\entete{blue}{Les événements et messages}

\begin{tcolorbox}[colback=yellow!10,colframe=yellow!50!black]
\textbf{Événement} : quelque chose qui se produit et déclenche un script.\\
\textbf{Message} : un signal envoyé par un lutin et reçu par un autre.
\end{tcolorbox}

\begin{tabularx}{\linewidth}{>{\bfseries}p{4.6cm}Y}
\toprule
Bloc & Rôle\\
\midrule
\texttt{quand drapeau cliqué} & Démarre le programme\\
\texttt{quand touche espace pressée} & Réagit au clavier\\
\texttt{quand ce lutin est cliqué} & Réagit au clic\\
\texttt{envoyer à tous salut} & Envoie un message\\
\texttt{quand je reçois salut} & Démarre un script à la réception\\
\texttt{envoyer à tous et attendre} & Attend la fin des scripts qui reçoivent\\
\bottomrule
\end{tabularx}

\begin{code}[Modèle]
Lutin 1 : quand ce lutin est cliqué $\rightarrow$ dire «~Bonjour~» $\rightarrow$ envoyer à tous salut\\
Lutin 2 : quand je reçois salut $\rightarrow$ dire «~Salut !~»
\end{code}

%==================================================================
\newpage
\section*{Annexe B : Fiche de pratique guidée}
\entete{green}{Pratique guidée : événements et messages}

\subsection*{3a. Exemple collectif : chat et chien}
Qui envoie ? \dotfill\quad Qui reçoit ? \dotfill\quad Nom du message : \dotfill

\subsection*{3b. En binôme : le jour et la nuit}
\begin{tabularx}{\linewidth}{>{\bfseries}p{3.6cm}Y Y}
\toprule
Élément & Mon script & Mon test\\
\midrule
Événement de départ & & \\
Message envoyé & & \\
Réception & & \\
Résultat visible & & \\
\bottomrule
\end{tabularx}

\begin{tcolorbox}[colback=yellow!10,colframe=yellow!50!black,title={Cartes d'aide à utiliser dans l'ordre},fonttitle=\bfseries\small]
\textbf{Aide 1 :} regarde le projet du 3a.\\
\textbf{Aide 2 :} Qui envoie ? Qui reçoit ? Quel nom donnes-tu au message ?\\
\textbf{Aide 3 :} demande à l'enseignant de créer le message avec toi.
\end{tcolorbox}

%==================================================================
\newpage
\section*{Annexe C : Niveau A (en dessous du niveau)}
\entete{orange}{Niveau A : je relie envoi et réception}

\subsection*{Activité 1 : relie chaque bloc à son rôle (les rôles sont mélangés)}
\begin{tabularx}{\linewidth}{Y Y}
\toprule
\textbf{Bloc} & \textbf{Rôle}\\
\midrule
\texttt{envoyer à tous salut} & Démarre le script quand on appuie sur une touche\\
\texttt{quand je reçois salut} & Le lutin envoie un signal\\
\texttt{quand touche espace pressée} & Le lutin démarre quand il reçoit le signal\\
\bottomrule
\end{tabularx}

\subsection*{Activité 2 : complète le dialogue}
Lutin 1 dit : «~\dotfill~» puis envoie le message \dotfill.\\
Lutin 2 reçoit le message et dit : «~\dotfill~».

\subsection*{Activité 3 : teste}
Écris ce qui se passe si on enlève le bloc \texttt{envoyer} : \dotfill

%==================================================================
\newpage
\section*{Annexe D : Niveau B (au niveau)}
\entete{blue}{Niveau B : une histoire en trois scènes}

\begin{tabularx}{\linewidth}{>{\bfseries}p{2.4cm}Y Y Y}
\toprule
Scène & Arrière-plan & Lutin et action & Message pour la suite\\
\midrule
1 & & & \\
2 & & & \\
3 & & & \\
\bottomrule
\end{tabularx}

\subsection*{Questions}
\begin{enumerate}
 \item Combien de messages as-tu utilisés ? \dotfill
 \item Que se passe-t-il si deux lutins reçoivent le même message ? \dotfill
\end{enumerate}

%==================================================================
\newpage
\section*{Annexe E : Niveau C (au-dessus du niveau)}
\entete{violet}{Niveau C : jeu à deux lutins}

\textbf{Mission :} crée un jeu dans lequel un lutin joueur évite un lutin ennemi.

\begin{tabularx}{\linewidth}{>{\bfseries}p{3.4cm}Y Y}
\toprule
Message & Qui l'envoie ? & Qui le reçoit et que fait-il ?\\
\midrule
\texttt{touche} & & \\
\texttt{fin} & & \\
\texttt{rejouer} & & \\
\bottomrule
\end{tabularx}

\textbf{Contraintes :} une variable \texttt{vies}, un message \texttt{fin} qui arrête les scripts, un bouton «~Rejouer~».

\textbf{Défi bonus :} ajoute un deuxième ennemi qui apparaît quand le score dépasse 10.

%==================================================================
\newpage
\section*{Annexe F : Corrections et attendus (enseignant)}

\subsection*{Annexe B}
3a : le chat envoie, le chien reçoit, message \texttt{salut}. 3b : événement \texttt{quand ce lutin est cliqué} ; envoi \texttt{nuit} ; réception \texttt{quand je reçois nuit} avec changement d'arrière-plan et réplique.

\subsection*{Annexe C}
\texttt{envoyer à tous salut} $\rightarrow$ «~Le lutin envoie un signal~» ; \texttt{quand je reçois salut} $\rightarrow$ «~Le lutin démarre quand il reçoit le signal~» ; \texttt{quand touche espace pressée} $\rightarrow$ «~Démarre le script quand on appuie sur une touche~».\\
Activité 3 : sans le bloc \texttt{envoyer}, le second lutin ne réagit pas.

\subsection*{Annexe D}
Au moins deux messages pour trois scènes. Deux lutins qui reçoivent le même message réagissent en même temps (en parallèle).

\subsection*{Annexe E}
\texttt{touche} : envoyé par l'ennemi quand il touche le joueur, reçu par le joueur qui perd une vie. \texttt{fin} : envoyé par le joueur à 0 vie, reçu par tous. \texttt{rejouer} : envoyé par le bouton.\\
\textbf{Critères de réussite :} messages clairement nommés, jeu qui se termine et peut redémarrer.

%==================================================================
\newpage
\section*{Annexe G : Exit ticket et grille d'observation}
\entete{gray}{Exit ticket : événements et messages}

\begin{enumerate}
 \item Quel bloc permet à un lutin de réagir quand un autre lutin lui envoie un signal ?\\ \lig
 \item Pourquoi est-il utile de donner un nom clair à un message ?\\ \lig
 \item \autoeval
\end{enumerate}
\textbf{Attendus :} 1. \texttt{quand je reçois\ldots}. 2. Pour retrouver facilement quel lutin doit réagir et éviter les erreurs.

\bigskip
\grille{Événement utilisé}{Message envoyé}{Message reçu}

\end{document}
